\chapter{Of Crimes and Their Punishments}
\label{art:crimes}

\articlenote{In which is set down how an offense is measured, how a first offense is forgiven, and how a citizen answers, in the mines or in the cleaning, for every offense that is not.}

\section{The Three Measures of Correction}

This Pact answers a broken law in one of three measures: a grace given for a first offense under Section 2; a term in the mines for an offense named under Section 3; or cleaning under Article 18 for an offense named among the gravest under Section 4.

\begin{clauses}
\item No offense is answered by more than one of these measures, and no citizen may be given a lesser measure for a gravest offense, nor a greater measure for an offense this Pact does not name among the gravest, whatever the Judge's own judgment of the citizen's character.
\item Where this Pact does not name the measure that answers an offense, the Judge shall rule under Article 6, and that ruling shall thereafter stand as though the measure were named here.
\end{clauses}

\section{The Grace of a First Offense}

A citizen's first offense under this Pact, save an offense named among the gravest under Section 4, is answered by a grace: a formal warning entered in the Judicial archive under Article 6, and no term in the mines.

\begin{clauses}
\item A grace is given once to a citizen in their lifetime, for the whole of the offenses this Pact does not name among the gravest, and not once for every kind of offense a citizen might commit.
\item A second offense, of any kind not among the gravest, is answered by the mines under Section 3, whatever grace the citizen was earlier given.
\end{clauses}

\section{Offenses Answered by the Mines}

The following, upon a second offense or upon a first offense the Judge finds too grave for grace, are answered by a term in the mines: the diversion or waste of food, feed, or stores under Article 11; the keeping of an animal without sanction under Article 11; standing responsible, under Article 2, for a birth without the lottery's grant of Article 3; speaking or teaching the forbidden wish under Article 16 Section 1; the obstruction of an officer under Article 16 Section 7, where that Section so provides; the willful and persistent neglect of assigned labor, referred by a Head of Department under Article 13 and found by Judicial to continue notwithstanding that referral; and the keeping, making, or use of any relic or restricted instrument or device forbidden under Article 8, Article 9, or Article 16, without license, save where Article 16 classifies that relic among the gravest.

\begin{clauses}
\item A term in the mines is a term of labor in the workings below the Down Deep, under the direction of the Head of Mines as Article 10 provides, for a span the Judge sets and the Judge alone may shorten.
\item A citizen serving a term in the mines remains a citizen for every purpose of this Pact, and returns, upon the term's ending, to the trade and floor Judicial assigns, which need not be the trade or floor held before.
\end{clauses}

\section{Offenses Among the Gravest}

The following, and no others save as this section provides, are named among the gravest offenses under this Pact, and are answered by cleaning under Article 18, without the grace of Section 2 or the hearing Article 6 grants for a lesser offense: tampering with, obstructing, or attempting private repair of the outward sensors, under Article 1; interference with the generator, the works of water and air, the labor of reclamation, or any system kept under Article 9, Section 7, whether by a citizen's own hand or by the keeping of a dwelling under Article 14; making a false attestation of a citizen's life or death under Article 2, or a false finding of a death's cause under Article 19, or a physician's failure to carry a suspected self-killing to Judicial as those Articles require; the taking, or the attempt to take, of a citizen's own life by any means not sought and granted under Article 18; the working of ore, or the directing of its working, beyond the limit Article 10 sets; the counterfeiting of a chit, the knowing keeping of a counterfeit, or the theft of or tampering with the silo's blank chits, under Article 12, or the counterfeiting of a seal-stamp under Article 21; a wrong done to another citizen with intent to kill or maim, under Section 5; and the keeping, making, or use of any relic or restricted instrument Article 16 classifies among the gravest.

\begin{clauses}
\item The Judge alone may name a further offense among the gravest, by a ruling under Article 6, and that ruling shall state plainly why the offense is found to threaten the whole silo and not merely the citizen who committed it.
\item Nothing named among the gravest offenses under this Article may be answered by any measure but cleaning, and no citizen sentenced under this section may be sentenced instead to the mines, whatever representation is made under Article 6.
\item Where a citizen's own life is taken and the citizen does not survive to be sentenced, Judicial's finding under this section is entered in the archive under Article 6 and closes the citizen's name under Article 2, and no further punishment is had; a citizen who survives an attempt under this section is sentenced to cleaning as for any other offense named in this Article.
\end{clauses}

\section{Offenses Against Another Citizen}

A citizen who does violence upon another, who takes what is not theirs, or who otherwise wrongs a citizen in body or property, is answerable under this Article, and the Judge shall set the measure --- grace, the mines, or, where the wrong was done with intent to kill or maim, among the gravest --- according to the wrong done and the harm that came of it.

\begin{clauses}
\item A wrong done under this section is Judicial's to judge, whether or not the wronged citizen brings it forward, but a citizen who brings a wrong forward and is found to have brought it falsely is answerable, under this same section, for the false bringing.
\item Where a wrong under this section, or an order given and later found unnecessary under Article 22, has cost a citizen property, labor, or standing, the Judge may order the offending party, or the silo itself where Article 22 so provides, to make restitution to the wronged citizen or the wronged citizen's household, in chits, goods, or labor, in addition to whatever measure of grace, the mines, or cleaning the wrong itself requires. Restitution is not itself a measure of correction under Section 1, and its ordering neither lessens nor replaces the measure that section requires.
\end{clauses}

\section{The Form of a Grace}

A grace is the gentlest measure this Pact gives, and it is given in a fixed form, that it be known for what it is.

\begin{clauses}
\item A grace is entered by the Clerk of Judicial under Article 6 in these terms and no others: the citizen's name and floor; the Article and Section offended, and the offense in a line; the day; and the words \emph{This is the grace of a first offense under Article 17, given once and not again}; the Judge signs the entry.
\item The grace is read to the citizen aloud by the Clerk, in the hall of Judicial or on the citizen's floor where the Judge so orders, and the citizen signs the entry, or the Clerk enters that the citizen would not; a copy is given to the citizen, and the citizen's household is told the grace was given, and no more.
\item A grace stands against a citizen's name for the whole of the citizen's life, and is read by the Judge before any later sentence under Section 2; it is struck only upon a ruling under Article 6 that no offense occurred, and for no other cause, and a grace struck is entered as struck and not removed.
\item A grace given is not a finding that the citizen is answerable for any other thing than the offense named, and no office may weigh a grace against a citizen in the assignment of a trade, a dwelling, a shadow, or a sanction; but the Head of a Department is told of a grace given to a citizen of the Department, for the Department's own keeping.
\item No grace is given for an offense a citizen committed in an office, as Section 9 provides, nor for the diversion of stores in an emergency under Article 22, nor where this Pact elsewhere provides that the mines answer a first offense.
\end{clauses}

\section{The Reckoning of a Term in the Mines}

A term is a span of seasons, counted from a day and ended on a day, and every day of it is entered.

\begin{clauses}
\item The Judge sets a term in whole seasons, not fewer than one and not more than twenty, and enters it with the sentence under Article 6; a term is not set in days or cycles, and a term of a part of a season is a term of the whole.
\item A term begins on the day the Head of Mines gives the receipt under Article 6, Section 12, and ends at the turn of the season on which the whole of its seasons have run; the Head of Mines enters the beginning and reports the ending under Article 10, and the Clerk enters both.
\item A term does not run on any day the citizen is unfit for the workings under Article 10, Section 8, and the days so lost are added at the term's end; a term does not run on the rest day either, the rest day being given the sentenced citizen as any other, but the rest day is counted in the seasons and is not added.
\item A term is shortened only by the Judge, upon the Judge's own reading of the Head of Mines' reports under Article 10, or upon a petition under Article 6 that the Head of Mines joins; no term is lengthened once set, whatever the citizen does below, an offense committed in the workings being answered as a further offense under this Article and entered as such.
\item A citizen returned from the mines under Article 6 who commits a further offense answered by the mines is sentenced to a term not less than the first, and a citizen sentenced to a third term is sentenced to twenty seasons, whatever the offense; a citizen who has served three terms is not sentenced to the mines again, and the Judge rules under Article 6, Section 2 what measure answers the offense instead.
\item The household of a citizen sentenced keeps its dwelling and its ration for the term as Article 3 provides, and the citizen's chits under Article 12 are the citizen's own still, entered to the ledger and not drawn until the term ends; the Chit-Keepers are told the term's beginning and its end.
\end{clauses}

\section{The Weighing Between Grace and the Mines}

Where an offense is not among the gravest, the Judge chooses between the grace and the mines, and the choosing is not left to the Judge's mood.

\begin{clauses}
\item In finding a first offense too grave for grace under Section 3, the Judge weighs, and enters the weighing of, these things and no others: the harm the offense did, to a citizen, to a floor, or to the silo; whether the citizen meant the harm, or meant the act and not the harm, or meant neither; whether the citizen offended alone or drew others in; whether the citizen offended in an office, as Section 9 provides; and whether the citizen came forward before the offense was found.
\item A first offense that did no harm, was not meant to, drew no other citizen in, and was not committed in an office is answered by the grace, and the Judge enters no weighing; a first offense in which any of those is otherwise is weighed, and the weighing is read to the citizen with the sentence.
\item In setting the span of a term under Section 7, the Judge weighs the same things and enters the same weighing, and sets the span nearer one season where fewer of them stand against the citizen and nearer twenty where more; a term set without a weighing entered is reduced by the Judge to one season upon the citizen's petition under Article 6.
\item A citizen who came forward before the offense was found, and who surrendered what the offense concerned, is given the grace where the offense is not among the gravest, whatever else the weighing shows, save an offense committed in an office; this is the same grace the Forgiveness Holiday under Article 4 gives for a relic, and stands whether or not a Holiday has been declared.
\end{clauses}

\section{The Offense of an Officer}

An office excuses no offense, as Article 1 provides; and an offense in an office is the graver for the office.

\begin{clauses}
\item A citizen who offends in the conduct of an office --- the Mayor, the Deputy Mayor, the Judge, the Sheriff, a Head of Department, or a citizen acting under an office's charge --- is given no grace for the offense, and a first offense so committed is answered by the mines where this Pact names that measure for it, and by the gravest where it names that; the office is vacated by the sentence, and the successor named as the office's own Article provides.
\item An offense of the Mayor is investigated by the Officers of Judicial and judged by the Judge as any citizen's, and the Mayor's orders under Article 4 do not bind Judicial in it; the Deputy Mayor acts under Article 5 from the sentence.
\item An offense of the Judge is investigated by the senior Officer of Judicial by first confirmation, who reports to the Mayor and to the Assembly under Article 20; the Mayor hears the matter as the Judge would hear it under Article 6, sets the measure this Pact names, and the Assembly removes the Judge under Article 20 upon the sentence; the Judge so sentenced is carried out of the office by the Officers as any citizen.
\item An offense of the Sheriff is investigated and judged as any citizen's, the Officers of Judicial acting in the Sheriff's place for the holding under Article 7; an offense of a Head of Department likewise, the Head's shadow or the senior citizen of the Department keeping the Department as acting Head until the Mayor confirms a successor.
\item An officer who offends is entered in the archive under the office as well as under the name, and the entry is read at the next sitting of the Assembly under Article 20 with the sentence, that the silo know the office was kept and not the officer.
\end{clauses}
